\section{Introduction and Vulnerability Overview}

\subsection{Vulnerability Overview}

This report investigates CVE-2026-31431, commonly referred to as
\emph{Copy Fail}, a local privilege escalation vulnerability in the Linux
kernel. The goal is to reproduce the vulnerability in an isolated laboratory
environment, understand the underlying bug, and use the observed behaviour to
develop and test detection and mitigation approaches.

The vulnerability affects \code{algif\_aead}, the Linux kernel's userspace
interface for AEAD algorithms exposed through \code{AF\_ALG} sockets. The
vulnerable code incorrectly performs cryptographic operations in place even
though the source and destination originate from different mappings. Combined
with \code{splice()}, this can cause page-cache-backed data from a readable file
to become part of a writable cryptographic scatterlist. Under the right
conditions, an unprivileged process can therefore modify the cached contents
of a file without having write permission to the file itself
\cite{ubuntu-copyfail-cve,linux-copyfail-fix}.

The published Copy Fail research turns this behaviour into a deterministic,
attacker-controlled four-byte write into the page cache of a readable file.
By repeatedly applying this primitive to a suitable setuid-root executable,
an otherwise unprivileged local user can modify the cached executable image
and ultimately execute code with root privileges
\cite{copyfail-research}.

Canonical classifies CVE-2026-31431 as a High-severity vulnerability with a
CVSS 3.1 score of 7.8. Besides local privilege escalation on the host,
Canonical also notes that the vulnerability may facilitate container escape
scenarios in affected environments
\cite{ubuntu-copyfail-cve}.

For this project, however, the interesting part is not simply whether the
published exploit can produce a root shell. The more useful question is what
actually happens on the system while the vulnerability is being exploited,
which parts of that behaviour can be observed, and whether the exploitation
path can be detected or restricted before the vendor patch is applied.

\subsection{Technical Background and Vulnerability Mechanism}

To understand Copy Fail, it is important to distinguish between a file on disk
and the copy of its contents currently held in memory. Linux uses the page
cache to keep recently accessed file data in memory, allowing subsequent reads
and executions to use the cached pages instead of repeatedly accessing the
underlying storage. A process with read access to a file may therefore cause
its contents to enter the page cache, but this should not give that process the
ability to modify those contents.

Copy Fail breaks this assumption through an unexpected interaction between the
page cache and the Linux kernel crypto API. Linux exposes parts of this API to
userspace through \code{AF\_ALG} sockets. For authenticated encryption,
\code{algif\_aead} provides the corresponding AEAD interface. AEAD
(Authenticated Encryption with Associated Data) combines encryption with
integrity protection and allows additional authenticated data to be included
without encrypting it.

Internally, cryptographic operations use scatterlists to describe the different
regions of memory involved in an operation. In the vulnerable implementation,
\code{algif\_aead} performs AEAD operations in place, meaning that the source
and destination of the cryptographic request ultimately share the same
scatterlist. This becomes problematic because the input and output do not
necessarily originate from the same type of memory
\cite{copyfail-research,linux-copyfail-fix}.

The second important part of the vulnerability is \code{splice()}. This system
call can move data between file descriptors without first copying it through a
conventional userspace buffer. When data from a readable file is spliced through
a pipe into an \code{AF\_ALG} socket, references to the file's page-cache pages
can become part of the crypto input scatterlist. During AEAD decryption,
\code{algif\_aead} copies the associated data and ciphertext into the receive
buffer, but chains the authentication-tag region from the original input onto
the destination scatterlist. As a result, file-backed page-cache memory can
become reachable through a scatterlist that the cryptographic operation treats
as writable \cite{copyfail-research}.

This alone does not yet produce the controlled write used by Copy Fail. The
final piece is the \code{authencesn} AEAD template. During decryption,
\code{authencesn} temporarily rearranges IPsec Extended Sequence Number data
and uses the destination scatterlist as scratch space. As part of this process,
it performs a four-byte write beyond the legitimate AEAD output region. Under
normal circumstances this affects the operation's destination buffer. Combined
with the vulnerable in-place layout in \code{algif\_aead}, however, that write
can instead reach one of the chained page-cache pages
\cite{copyfail-research}.

The attacker can control both the position of this write and the four-byte
value being written. Even though the subsequent authentication check fails,
the modification is not rolled back. The affected page is also not marked
dirty for normal filesystem writeback, meaning the file on disk remains
unchanged while its cached in-memory representation has been modified. The
result is therefore a deterministic four-byte write into the page cache of a
file that the attacker only needs permission to read
\cite{copyfail-research}.

This limited primitive is the central building block of Copy Fail. It does not
grant root privileges by itself, but it creates something much more interesting
from an exploitation perspective: the ability to modify executable data in
memory without having normal write access to the underlying file. How this
primitive is turned into privilege escalation is discussed in the following
section.

\subsection{Security Impact}

At first glance, a controlled four-byte write may appear too limited to be useful
for privilege escalation. The important limitation, however, is only the amount
of data written during a single operation. Because the primitive is deterministic
and repeatable, an attacker can perform multiple writes at controlled offsets and
gradually modify a larger region of a cached file
\cite{copyfail-research}.

The published Copy Fail exploit demonstrates this by targeting
\code{/usr/bin/su}, a setuid-root executable commonly available on Linux systems.
Successive four-byte writes are used to modify executable code in the cached
image of the binary. The file itself remains owned by root and retains its setuid
permissions because neither its filesystem metadata nor its persistent contents
need to be changed.

When the modified executable is started, Linux uses the poisoned page-cache
contents while still applying the privileges associated with the original
setuid-root file. The injected code therefore executes in a privileged context,
allowing a local unprivileged user to obtain root privileges. In other words, the
attacker never needs legitimate write access to the executable. Read access,
combined with the vulnerable kernel path, is enough to cross the privilege
boundary \cite{copyfail-research}.

An additional characteristic of Copy Fail is that the corrupted page is not
marked dirty for normal filesystem writeback. The modification therefore exists
in memory while the corresponding data on persistent storage remains unchanged.
This makes the corruption temporary from the filesystem's perspective, but that
does not reduce its immediate impact: the attacker only needs the modified
executable to run once to turn the page-cache write into a root shell
\cite{copyfail-research}.

The impact is also not necessarily limited to processes running directly on the
host. Canonical notes that the same vulnerability may facilitate container escape
scenarios in affected environments, extending the security implications beyond
a conventional local privilege escalation
\cite{ubuntu-copyfail-cve}.

The significance of Copy Fail is therefore much larger than the size of the
individual write suggests. A four-byte primitive is enough to undermine the
integrity of privileged executable code and ultimately turn an unprivileged local
user into root.

\subsection{Vendor Response and Mitigation}

Following the public disclosure of Copy Fail, Canonical initially provided a
temporary mitigation that removed access to the vulnerable kernel path rather
than attempting to correct the underlying bug. An update to the \code{kmod}
package prevented the \code{algif\_aead} module from being loaded, making the
affected AF\_ALG AEAD interface unavailable through that module
\cite{ubuntu-copyfail-cve}.

If the module was already loaded, it additionally had to be unloaded or the
system rebooted for the mitigation to become effective. Canonical explicitly
treated this as a temporary measure. Applications that rely on the kernel AEAD
interface may fall back to userspace cryptographic implementations, but this
behaviour is not guaranteed and disabling the module may therefore cause
compatibility or functionality problems
\cite{ubuntu-copyfail-cve}.

The original Copy Fail research also proposes restricting the attack surface by
preventing creation of AF\_ALG sockets through mechanisms such as seccomp. Like
disabling \code{algif\_aead}, this can prevent the vulnerable path from being
reached, but it does not correct the vulnerability itself
\cite{copyfail-research}.

The central upstream fix instead addresses the underlying design problem in
\code{algif\_aead}. The kernel change titled
\emph{Revert to operating out-of-place} removes the problematic in-place
operation introduced by an earlier optimization. Because the source and
destination originate from different mappings, the kernel developers concluded
that there was no practical benefit in sharing the same buffers. Restoring
out-of-place operation keeps the cryptographic destination separate from the
file-backed input pages and removes the condition that allowed page-cache memory
to become writable through the AEAD operation
\cite{linux-copyfail-fix}.

Distribution kernel updates include this change together with related fixes in
the AF\_ALG and authentication code. Once a corrected kernel is installed and
the system has been rebooted into it, the temporary module-level mitigation is
no longer required. Canonical therefore recommends updating to a fixed kernel
instead of relying permanently on disabling \code{algif\_aead}
\cite{ubuntu-copyfail-cve}.

This distinction is important for the experimental part of this project.
Disabling the vulnerable interface can stop exploitation, but it does so by
removing functionality. The kernel patch instead removes the vulnerable
behaviour itself. Both approaches provide useful reference points when evaluating
whether additional detection or mitigation techniques can identify or interrupt
Copy Fail without simply disabling the affected subsystem.

\subsection{Research Objectives and Questions}

The objective of this project is not limited to successfully executing the
published Copy Fail proof of concept. Reproducing the privilege escalation is
treated as the starting point for a broader investigation into how the
vulnerability behaves, how that behaviour can be observed, and which defensive
measures can meaningfully interrupt the exploitation path.

The experimental workflow therefore follows the vulnerability through several
states: first understanding and reproducing it, then observing and detecting its
behaviour, applying a compensating control, comparing that control with the
vendor mitigation, and finally retesting against a patched kernel.

The investigation is guided by the following research questions:

\begin{description}

    \item[RQ1 Reproduction.]
    Can Copy Fail be reproduced reliably in the isolated laboratory environment
    by an ordinary, unprivileged local user?

    \item[RQ2 Observability.]
    Which system-level behaviours can be observed consistently while the
    vulnerability is being exploited, and which of these provide useful
    detection opportunities?

    \item[RQ3 Detection.]
    Can the exploitation technique be detected using behavioural indicators
    rather than artifacts that are specific to a particular proof of concept,
    such as its filename or executable hash?

    \item[RQ4 Compensating Control.]
    Can a defensive control interrupt the known exploitation path before the
    kernel is patched, and can it do so without unnecessarily removing benign
    functionality?

    \item[RQ5 Vendor Comparison.]
    How does the custom compensating control compare with Canonical's temporary
    mitigation, which prevents \code{algif\_aead} from being loaded?

    \item[RQ6 Patch Validation.]
    Does reproduction fail after installation of a corrected kernel, and do the
    observations from the final retest match the expected behaviour of the
    upstream fix?

\end{description}

Together, these questions turn the published exploit into an experimental
starting point rather than the final result. A successful root shell proves that
the vulnerability can be reproduced, but the more useful outcome is understanding
what distinguishes the vulnerable behaviour, how defenders can react to it, and
what changes once the underlying kernel flaw is actually removed.

\subsection{Scope and Safety}

All experimentation in this project is restricted to systems owned and
controlled for the purpose of the research and placed inside an isolated
laboratory environment. No third-party or internet-facing systems are included
in the test scope. The deliberately vulnerable reproduction target is kept
separate from the patched reference system so that a known-good comparison state
remains available throughout the project.

Reproduction is performed from a dedicated unprivileged local account. For the
purpose of this project, privilege escalation is considered successfully
demonstrated once that account can execute a process with UID 0 and the result
can be recorded using a minimal command such as \code{id}. No persistence,
credential theft, lateral movement, or unrelated post-exploitation activity is
required or performed.

The experimental scope is limited to the published Copy Fail exploitation path
involving the AF\_ALG AEAD interface and modification of file-backed page-cache
contents. Potential container escape is relevant to the vulnerability's overall
impact, but it is not reproduced as part of this project. Other Linux kernel
vulnerabilities and unrelated privilege-escalation techniques are likewise
outside the scope of the experiment.

System state is recorded before vulnerability-specific changes or exploit
execution. Snapshots and phase checkpoints are used so that an ambiguous or
modified system can be restored rather than treated as a clean reference.
Experimental evidence is collected separately from the analytical report. Raw
artifacts are preserved privately with integrity verification, while material
intended for public release is sanitized where necessary and receives separate
integrity verification after sanitization.

These boundaries keep the experiment focused on understanding, observing, and
defending against Copy Fail while ensuring that the reproduction remains
controlled, reversible, and limited to the behaviour required to answer the
research questions.